\SourceLeaf{531}
\begin{center}
\begin{tabular}{rrr@{\hspace{2em}}rrr}
Chilias & \textsc{Gauss} & Wahrer Werth & Chilias & \textsc{Gauss} & Wahrer Werth \\
$20$ & $102$ & $104$ & $546$ & $68$ & $69$ \\
$159$ & $87$ & $77$ & $601$ & $75$ & $76$ \\
$199$ & $96$ & $86$ & $625$ & $68$ & $78$ \\
$206$ & $85$ & $83$ & $668$ & $73$ & $74$ \\
$245$ & $78$ & $88$ & $675$ & $69$ & $73$ \\
$289$ & $85$ & $77$ & $784$ & $74$ & $75$ \\
$290$ & $84$ & $85$ & $800$ & $81$ & $71$ \\
$334$ & $80$ & $81$ & $879$ & $68$ & $78$ \\
$352$ & $80$ & $81$ & $985$ & $74$ & $70$ \\
$501$ & $78$ & $79$ &  &  &  \\
\end{tabular}
\end{center}
\par Die in dem Briefe von \textsc{Gauss} an \textsc{Encke} erwähnte Formel \textsc{Encke}’s scheint die folgende
\[
\frac{n}{\log n}\overset{2\log n}{\surd}10
\]
zu sein, welche \textsc{Encke} in einem Briefe an \textsc{Gauss} vom 4. Dec. 1849 mittheilt.
\par Die \textit{Tafel der Anzahl der Classen binärer quadratischer Formen} gibt die Anzahl der Genera und Classen so wie den Index der Irregularität für die negativen Determinanten in den Hunderten 1 bis 30, 43, 51, 61, 62, 63, 91 bis 100, 117 bis 120, dann noch in einer besondern Zusammenstellung für die des 1. 3. und 10\textsuperscript{ten} Tausend, für die 800 ersten von der Form $-(15n+7)$ und $-(15n+13)$, sowie für einige sehr grosse Determinanten, ferner für die positiven Determinanten des 1. 2. 3. 9. 10\textsuperscript{ten} Hundert und für einige andere. Die Handschrift besteht aus einzelnen Zetteln, auf denen die Tafeln verschiedenartig eingerichtet sind, z. B. ist bei den ältern das Wort Ordo statt Genus gebraucht, so bei den einzelnen Centaden mit Ausnahme der 9. und 10. positiver Determinanten, dann aber auch bei einzelnen vorläufigen Zusammenstellungen in Chiliaden.
Zur leichtern Uebersicht ist hier überall die Bezeichnung der Disqu. Arithm. gewählt, auch die grössten und kleinsten Quotienten aus der Anzahl der Classen dividirt durch den Determinanten, sowie die Anzahl der Determinanten, für welche der Quotient innerhalb gewisser Grenzen fällt, sind wegen Mangel an Raum nicht unter die einzelnen Centaden gesetzt sondern am Ende der Tafel für die negativen Determinanten zusammengestellt. Aus einigen übrig gebliebenen Aufzeichnungen scheint hervorzugehen, dass \textsc{Gauss} zuerst die Classen für die Determinanten berechnet hat, die demselben Hundert und demselben Reste bei dem Theiler 15 angehören. Die Determinanten dieser Abtheilungen sind dann nach der Anzahl der Genera und Classen und zuletzt alle die demselben Hundert angehörigen auf die hier wiedergegebene Weise geordnet.
Den Tafeln der einzelnen Centaden sind manche spätere Berichtigungen eingefügt, nicht aber den Zusammenstellungen in Tausenden. Zeitbestimmungen enthalten nur die beiden Tafeln mit den Determinanten der Form $-(15n+7)$ und $-(15n+13)$ nemlich resp. ‘\textit{Expl. In. Febr. 1801}’ und ‘\textit{Expl. 27 Febr. 1807.}’
\par In diesen Tafeln habe ich unter andern die folgenden Fehler bemerkt, denen ich hier zur leichtern Controle die Periodenzahlen der Fundamentalclassen wie z. B. $4\pdot\ 4\pdot\ 2$ bei dem Determinanten
